%---------------------------------------------------------
\subsection{Vista C -- Analista delle prestazioni e del gemello digitale}
\label{sec:req-C}
L'analista lavora come gestore di progetto relativo ad un singolo committente di System Logistics e segue da remoto il rendimento
dell'impianto dopo l'avviamento~\cite{sl_cs}. Deve capire se le \concetto{macchine} lavorano come
dovrebbero e, per farlo, mette a confronto i dati raccolti sul campo con le previsioni del
\concetto{gemello digitale} (la copia virtuale dell'impianto su cui si fanno le
simulazioni). Consulta le misure dei \concetto{sensori} e gli \concetto{eventi} di
ogni \concetto{macchina}, controlla gli \concetto{indicatori} calcolati a fine turno e ricalibra il
\concetto{gemello digitale} quando le sue previsioni si allontanano dalla realtà. Prepara inoltre gli
\concetto{scenari} da simulare, confronta valori reali e previsti e segnala al tecnico dell'assistenza
le \concetto{macchine} che rendono meno del previsto. Le sue simulazioni servono anche al committente,
che prima di un periodo di picco vuole sapere se l'impianto reggerà.

\paragraph{Sensori e misure.}
Gran parte dei dati di campo arriva dai \concetto{sensori}. Di ognuno si registrano un codice, la
grandezza misurata (per esempio la corrente assorbita da un motore o una temperatura), l'unità di misura
e ogni quanto viene letto. \evid{G}{Un \concetto{sensore} può essere di macchina, se è montato su una
\concetto{macchina}, oppure ambientale}: in questo caso controlla una zona dell'impianto e ha un
intervallo di valori ammessi, come i termometri dell'area di stoccaggio, dove le bevande non devono
scaldarsi troppo. Ogni lettura di un \concetto{sensore} produce una misura, con istante e valore.

\paragraph{Eventi.}
Oltre alle misure, le \concetto{macchine} comunicano degli \concetto{eventi}: li invia il controllore di
bordo (PLC) attraverso SYSTORE, ognuno con il proprio istante. \evid{G}{Gli \concetto{eventi} si dividono in cambi di stato, allarmi e fine missione}. Un
cambio di stato registra lo stato lasciato e quello nuovo (in lavoro, in attesa, in errore, in
manutenzione). Un allarme ha un codice e una gravità, che va dal semplice avviso fino al blocco della
\concetto{macchina}. Un \concetto{evento} di fine missione indica la missione che la
\concetto{macchina} ha appena concluso.

\paragraph{Indicatori.}
Gli \concetto{indicatori} di prestazione (KPI) sono definiti una volta sola. Ciascuno con
codice, nome, formula di calcolo e unità di misura. Si usano la disponibilità, il throughput (UDC
movimentate all'ora), il tempo medio di ciclo e l'energia consumata per missione. Alla fine di ogni
turno, individuato da data e fascia oraria, il sistema calcola il valore reale di ogni
\concetto{indicatore} per ogni \concetto{macchina}: la disponibilità deriva dalla durata degli stati,
throughput e tempo di ciclo dalle missioni concluse, l'energia dalle misure dei \concetto{sensori}. Il
termine di paragone è il valore previsto dal \concetto{gemello digitale}: ogni \concetto{indicatore} ha
una soglia che dice di quanto il valore reale può allontanarsene prima che la \concetto{macchina} venga
segnalata. Se per esempio la soglia della disponibilità è di tre punti e la previsione è del 96\%, un
valore reale del 92\% fa scattare la segnalazione. Alcuni \concetto{indicatori}, come la disponibilità,
hanno anche un valore obiettivo fissato nel contratto di assistenza.

\paragraph{Gemello digitale.}
Il \concetto{gemello digitale} riproduce in simulazione le \concetto{macchine} dell'impianto e il loro
funzionamento. Ogni volta che viene ricalibrato sui dati di campo ne nasce una nuova versione, con numero
progressivo, data di calibrazione ed errore medio delle previsioni rispetto ai valori reali.
\evid{PO}{Ogni versione del \concetto{gemello digitale} è composta dalle repliche virtuali delle singole
\concetto{macchine}}: una replica corrisponde a una \concetto{macchina} e ne contiene i parametri di
funzionamento (velocità, accelerazioni, tempi di presa e deposito), che possono cambiare da una versione
all'altra. Per ogni versione si indica anche il livello di integrazione con
l'impianto~\cite{kritzinger2018}: modello digitale se i dati sono scambiati a mano, ombra digitale se
arrivano in automatico dall'impianto alla simulazione, gemello vero e proprio se viaggiano in automatico
nei due sensi.

\paragraph{Scenari e simulazioni.}
Uno \concetto{scenario} stabilisce le condizioni in cui far lavorare il \concetto{gemello digitale}. Ha
un nome, una descrizione e le condizioni ipotizzate, come un diverso profilo degli ordini o la modifica
di alcune repliche, per simulare il fermo di una \concetto{macchina} o una velocità maggiore. Lo
\concetto{scenario} di riferimento riproduce le condizioni effettive di ogni turno e serve al confronto
quotidiano: dagli scarti tra le sue previsioni e i valori reali si ricava anche l'errore medio di ogni
versione. Gli altri \concetto{scenari} sono ipotesi del tipo «cosa succede se», come un picco di ordini
prima dell'estate o l'aggiunta di un turno notturno. \evid{IO}{Ogni simulazione è un'esecuzione
concreta di uno \concetto{scenario}} su una versione del \concetto{gemello digitale}, con data e durata,
e restituisce il valore previsto di ogni \concetto{indicatore} per ogni \concetto{macchina}. La
simulazione resta legata alla versione usata, così si possono confrontare previsioni ottenute con
calibrazioni diverse.
